\subsection{The pace, and where the cost went}

An account of this method that did not state its speed would be evasive, since
speed is the whole of what changed. The figures are in the repository's own
history and can be recovered from it.

The transcription of 667 sheets took nine days, from 5 to 13 September 2026, in
sixty-two passes of twelve sheets. Nobody keyed a character: a pass is a
photograph read and a file written, and the twelve sheets of a batch take the
model minutes. In the same period the repository received 210 commits across
thirteen days --- the scripts that derive the accounts, the formats, the
materials and the calendar (10,900 lines), the site that places the reading
beside the photograph (15,500 lines), the TEI converter and its validator, and
later the 1,800 lines of Lean that certify the arithmetic. A transcription
project of this size, keyed and proofed by people, is a year of somebody's
salary; this was a fortnight of one person's attention and a model's time.

That comparison is the easy half, and it is not the interesting one. What the
method actually changes is the second operation, the one that has no name in
the traditional workflow: \emph{restructuring}. Once the transcription is a
constrained language rather than a document, every question one might put to
the archive becomes a script of a few hundred lines --- what was paid and by
whom, which formats recur, when the colourman changed, on what days the books
were written in. Each of those is a census this edition publishes, each was
written in an afternoon, and each is regenerated on every build. The cost of
asking a new question of the corpus has fallen by more than the cost of reading
it, because reading is done once and asking is done for ever.

The cost did not disappear; it moved, and it moved to the two places this
article has been describing. The first is the definition of rules. The Book~IV
reader changed its rule four times across the volume's fifty-four years, and
every one of those changes is an editorial decision about what the writer meant
that no model and no script could take: whether a bare figure is a subtotal,
whether a « less » line is a deduction, whether three payments are three sales
or one. Writing them down --- in the method note, then as Lean functions proved
for every input --- took longer than transcribing the volume. The second is
verification. 667 sheets are read and none is checked. A pass costs minutes and
a human check costs an hour a batch, so the reading has produced, at the pace
described above, a debt of some fifty hours that the method cannot discharge
and does not pretend to.

The same arithmetic applies to the sibling corpora, and states the size of what
is left. The Grothendieck fonds stands at 1,661 pages in 176 passes, nine pages
to a pass, against sixteen thousand openly accessible pages at Montpellier: at
that rate the reading is some 1,700 further passes, which is months of machine
time and, at an hour a batch, about two thousand hours of human checking ---
that is, a reading a machine can finish this year and a verification no
individual can. The mathematicians' manuscripts in Gallica stand at 548 views
in 34 passes, sixteen views to a pass, against 25,258 views cataloged: some
1,500 further passes. Neither number is an argument for doing it; both are the
argument of this article, which is that the transcription is no longer the
bottleneck and the checking is, and that an edition which cannot say which of
its sheets a person has read is publishing a machine's confidence as though it
were scholarship.
